% Conclusiones RELEVANTES del estudio realizado. Deber\'ia responder a las preguntas: ?`todas las t\'ecnicas resuelven el mismo problema o hay algunas diferencias?, ?`En qu\'e se parecen o difieren las t\'ecnicas en el contexto del problema?, ?`qu\'e limitaciones tienen?, ?`qu\'e t\'ecnicas o estrategias son las m\'as prometedoras?, ?`existe trabajo futuro por realizar?, ?`qu\'e ideas usted propone como lineamientos para continuar con investigaciones futuras?

Se pudo observar que para el problema de TDPDP se ha realizado tanto enfoques de técnicas completas como de heuriticas y metaheuristicas. En un $53\%$ de los casos se utilizaron tecnicas heuristicas y metaheuristicas para resolver los problemas, obteniendo soluciones iniciales obtenidas con métodos completos. Esto fue debido a la naturaleza NP-Hard del problema, la cual fue evidenciada dentro del desarrollo del informe.

En base a las soluciones se evidencio la diferencia en el enfoque de como resolver el problema. Dado que el problema es un VRP modificado, se observo que una de las técnicas más utilizadas fue la de resolver primeramente un viaje para el camión y luego ir revisando en que secciones se podria cambiar la entrega por la de un dron. También resulto util diferenciar los problemas según cantidad de drones y camiones, para así enfocarse en solucionar una parte del problema que sea menos compleja que el total.

Las mayorias de limitaciones se pueden resumir a la complejidad del problema. Dado que es NP-Hard, se tarda mucho tiempo en obtener una solución del problema la cual no se vea afectada por largos tiempos de computación. Esto es debido al tamaño del espacio de búsqueda y las restricciones de sincronizar un dron con un camión. Un factor que se vio fue como las restricciones van aumentando mientras el problema se vuelve mas cercano a la realidad. En los primeros modelos, las restricciones eran muy ideales y no consideraban tanto la realidad del problema. Trabajos posteriores arreglan esto, agregando gasto de bateria variable según el peso del dron, o permitiendo que el dron se una a mitad de camino con un camión, etc.

Las tendencias muestran que el trabajo se encuentra focalizado en hacer uso de técnicas mixtas, que combinan el modelo matemático puro con las heuristicas para la resolución de los nodos utilizado spara el dron. El trabajo futuro que queda por realizar seria hacer un análisis más profundo de las diferencias entre los modelos (según cantidad de drones, camiones, batería) y como estos modelos se relacionan con la realidad de la logistica, para que así este trabajo pueda desarrollar mayor valor en empresas de entrega.
